% New executed retrospective public EEG checks.
\subsection{Alternative preprocessing and fixed-classifier sensitivity}
\label{sec:p2-real-extra-robustness}
We evaluated two retrospective robustness questions after the
primary content-overlap finding. Neither is a new clinical
cohort, independently registered replication or causal
estimate. First, exact whole-record identity grouping
of the same 50 publicly released Liu2024 IDs was
repeated using a different processing configuration
(6--35 Hz filtering and 128 Hz resampling rather
than 8--30 Hz/160 Hz), without reading classifier
outcomes. Released IDs 4 and 5 again contributed
identical complete processed 40-trial and event-label
arrays. The cohort again contained 49 whole-record
content groups, with an array shape of
$40\times29\times512$ per released ID.
This is one preprocessing sensitivity analysis,
not an independent raw EDF reader.
The executed record is archived at
\url{https://github.com/lindicaphxag-tech/kaggle/actions/runs/38319497449}.

Second, we refitted two fixed-hyperparameter
tangent-space classifiers on public EEG, leaving
out each of the same 49 held recording contents.
For each held content, the original model used
the 49 other nominal released IDs, including
any overlapping counterpart; the purged source
model used the other 48 content representatives.
Each fit independently reestimated its source
Riemannian tangent reference, standardisation
and classifier weights. The held target labels
were used only for final ROC-AUC calculation,
not for source fitting or parameter selection.
The compared classifiers were logistic regression
with $C=1$ (liblinear) and ridge classification
with $\alpha=10$. Their training losses differ
but both consume the same kind of covariance
tangent features.

\begin{table}[tb]
\centering\small
\caption{New retrospective model-family sensitivity
on the same 49 held digital recording-content groups.
All entries are purged-source-fit minus nominal
source-fit ROC-AUC. The 49 fitted source training
sets overlap extensively, so these rows are
not 49 independent model replications.}
\label{tab:p2-multimodel}
\setlength{\tabcolsep}{4pt}
\begin{tabular}{p{.31\linewidth}rrrr}
\toprule
Frozen model & All 49 mean & ID 4 & Other 48 & Median \\
\midrule
Logistic $C=1$
& $-0.004796$ & $-0.212500$ & $-0.000469$ & $0.000000$\\
Ridge $\alpha=10$
& $-0.003469$ & $-0.222500$ & $+0.001094$ & $+0.002500$\\
\bottomrule
\end{tabular}
\end{table}

For logistic regression, the original and purged
means were $0.580867$ and $0.576071$,
respectively. The corresponding ridge means were
$0.582296$ and $0.578827$. Logistic regression
had 22 negative and 21 positive paired changes
and ridge had 18 negative and 28 positive ones.
The pronounced ID 4 content-group change
therefore persists under two additional fixed
training objectives, while the remaining
48 groups show no consistent negative mean shift.
The new experiments support
\emph{model-specification sensitivity} and a
local overlap-related effect, not a universal
AUC correction for previously unseen patients.
The reproducible, per-record numerical results
are in the successful workflow artifact
\url{https://github.com/lindicaphxag-tech/kaggle/actions/runs/38319072438}.
